\documentclass[a0paper,landscape]{ImperialPoster}

% Essential packages
\usepackage{amsmath,amssymb}
\usepackage{graphicx}
\usepackage{multicol}
\usepackage{booktabs}
\usepackage{enumitem}
\usepackage{adjustbox}

% Compact list formatting
\setlist{nosep,leftmargin=12pt,labelsep=8pt}

% Define color scheme  
\definecolor{imperialblue}{RGB}{0,62,116}
\definecolor{imperialred}{RGB}{214,0,28}

% Custom commands
\newcommand{\alertignored}{80\%}

% Reduce paragraph spacing
\setlength{\parskip}{18pt}

% Title information
\postertitle{Adaptive Cyber Defense Through Reinforcement Learning: Dynamic Decoy Deployment Against Advanced Persistent Threats}
\posterauthors{Mo Alkhazaleh \\ Department of Computing, Imperial College London}

\begin{document}
\titlesection

\begin{multicols}{4}
\raggedcolumns

% COLUMN 1 - Problem & Method Overview
\section{The Cyber Challenge}

\textbf{Asymmetric Advantage:} Attackers need one vulnerability; defenders must protect everything

\textbf{Alert Fatigue:} \alertignored{} of alerts ignored due to volume  

\textbf{Static Defense Limitations:} Traditional approaches remain static while threats evolve

\vspace{12pt}
\begin{center}
\includegraphics[width=0.9\columnwidth]{ADAPTIVE_BLUE_DEFENSE_PRESENTATION/slide2_problem_solution.png}
\end{center}

\section{Our Approach}

\textbf{Reinforcement Learning Solution:} Dynamic decoy deployment that adapts to attack patterns

\begin{itemize}
\item Continuous learning from network interactions
\item Strategic decoy placement optimization  
\item Real-time adaptation to evolving threats
\item Resource-efficient defensive strategies
\end{itemize}

\vspace{12pt}
\begin{center}
\includegraphics[width=0.9\columnwidth]{ADAPTIVE_BLUE_DEFENSE_PRESENTATION/slide3_reinforcement_learning.png}
\end{center}

% COLUMN BREAK
\columnbreak

% COLUMN 2 - Technical Method
\section{PPO-Based Defense}

\textbf{Proximal Policy Optimization} provides stable learning for defensive agents:

\begin{itemize}
\item Clipped surrogate objective prevents destructive updates
\item Value function learning reduces gradient variance
\item Sample-efficient learning from limited interactions
\end{itemize}

\vspace{12pt}
\begin{center}
\includegraphics[width=0.95\columnwidth]{ADAPTIVE_BLUE_DEFENSE_PRESENTATION/slide8_ppo_algorithm.png}
\end{center}

\section{Network Environment}

\textbf{Cyberwheel Simulation:} Realistic enterprise network environment

\vspace{12pt}
\begin{center}
\includegraphics[width=0.95\columnwidth]{ADAPTIVE_BLUE_DEFENSE_PRESENTATION/slide5_network_environment.png}
\end{center}

\textbf{Reward Structure:}
$$R_t = 10v_r\mathbf{1}_{decoy} - v_r\mathbf{1}_{asset} - 2\mathbf{1}_{deploy} - \mathbf{1}_{remove}$$

\begin{itemize}
\item High reward for decoy interactions (successful deception)
\item Penalty for real asset compromise  
\item Action costs encourage efficient resource use
\end{itemize}

% COLUMN BREAK
\columnbreak

% COLUMN 3 - Results
\section{Experimental Results}

\textbf{Performance Comparison:}

\vspace{8pt}
\begin{center}
\begin{tabular}{lcc}
\toprule
\textbf{Method} & \textbf{Detection Rate} & \textbf{False Positives} \\
\midrule
Random Placement & 45\% & 32\% \\
Static Optimal & 62\% & 28\% \\
Heuristic Rules & 58\% & 25\% \\
\textbf{PPO Agent} & \textbf{78\%} & \textbf{15\%} \\
\bottomrule
\end{tabular}
\end{center}

\vspace{12pt}
\textbf{Learning Convergence:}

\begin{center}
\includegraphics[width=0.95\columnwidth]{Training_Convergence.png}
\end{center}

\textbf{Key Performance Metrics:}
\begin{itemize}
\item 73\% improvement in detection rate over random
\item 53\% reduction in false positive alerts
\item Stable convergence within 500k training steps
\item Consistent performance across network sizes
\end{itemize}

\section{Strategic Behaviors}

The PPO agent learns sophisticated defensive strategies:

\begin{itemize}
\item \textbf{Predictive Placement:} Anticipates attacker movement patterns
\item \textbf{Resource Optimization:} Balances coverage with computational cost
\item \textbf{Adaptive Response:} Adjusts strategy based on attack sophistication
\item \textbf{Vulnerability Prioritization:} Focuses protection on high-value assets
\end{itemize}

% COLUMN BREAK
\columnbreak

% COLUMN 4 - Impact & Future Work
\section{Key Contributions}

\textbf{1. Novel RL Framework}
First application of PPO to dynamic decoy deployment in enterprise networks

\textbf{2. Adaptive Defense Paradigm}  
System learns and counters evolving attack strategies in real-time

\textbf{3. Empirical Validation}
Comprehensive evaluation against multiple threat models and baselines

\textbf{4. Scalable Implementation}
Demonstrated effectiveness on networks up to 100 nodes

\vspace{12pt}
\begin{center}
\includegraphics[width=0.9\columnwidth]{ADAPTIVE_BLUE_DEFENSE_PRESENTATION/slide11_contributions.png}
\end{center}

\section{Research Impact}

\textbf{Paradigm Shift:} From static to adaptive defense

\textbf{Practical Benefits:}
\begin{itemize}
\item Significantly outperforms static approaches
\item Reduces operational burden on security teams  
\item Provides proactive rather than reactive security
\item Scales to enterprise network complexity
\end{itemize}

\section{Future Directions}

\begin{center}
\includegraphics[width=0.9\columnwidth]{ADAPTIVE_BLUE_DEFENSE_PRESENTATION/slide13_future_directions.png}
\end{center}

\textbf{Next Steps:}
\begin{itemize}
\item Multi-agent coordination across network segments
\item Adversarial co-evolution with offensive agents
\item Real-world deployment and validation
\item Integration with existing security infrastructure
\item Explainable AI for analyst decision support
\end{itemize}

\textbf{Long-term Vision:} Intelligent defensive systems that continuously adapt and evolve, maintaining effectiveness against sophisticated adversaries while minimizing false alarms.

\end{multicols}

\end{document}